\begin{abstract}
Báo cáo này trình bày tổng quan về phương pháp \tr{kernel} (kernel methods),
một trong những công cụ nền tảng và mạnh mẽ nhất trong học máy giúp giải quyết
các bài toán phân tách phi tuyến. Bằng cách sử dụng thủ thuật \tr{kernel}
(kernel trick), dữ liệu được ánh xạ ngầm định sang các không gian đặc trưng
nhiều chiều mà không đòi hỏi chi phí tính toán trực tiếp khổng lồ. Cụ thể, báo
cáo tập trung làm rõ ba nội dung cốt lõi: thứ nhất, cơ sở lý thuyết về
\tr{kernel} đối xứng xác định dương (PDS kernels) và sự liên hệ chặt chẽ với
không gian Hilbert \tr{kernel} tái tạo (RKHS); thứ hai, kỹ thuật thiết kế
\tr{kernel} chuỗi (sequence kernels) để xử lý các loại dữ liệu rời rạc, phi cấu
trúc như chuỗi ký tự; và cuối cùng là phương pháp xấp xỉ ánh xạ \tr{kernel} đặc
trưng bằng đặc trưng Fourier ngẫu nhiên, nhằm giải quyết bài toán nút thắt cổ
chai về bộ nhớ và thời gian tính toán trên các tập dữ liệu quy mô lớn.  Thông
qua đó, báo cáo cung cấp một góc nhìn toàn diện từ cơ sở toán học chặt chẽ đến
các giải pháp tối ưu hóa khả năng mở rộng trong thực tiễn.
\end{abstract}
